%! Inference and Experiments — Unit 3, Lesson 3.5 — Lesson Plan
% GRADUAL-RELEASE plan (warm-up / guided notes and practice / debrief / close and start the
% homework) at the COMPACT budget: the whole lesson is finished in one 55-minute period.
% THERE IS NO GROUP ACTIVITY IN THIS COURSE and NO INDEPENDENT PRACTICE SET in the notes — the
% homework is the individual practice, started in class. Teacher-facing. Every box is `skillbox`
% (breakable) with a \boxguard ahead of it; no `fixedskillbox`, no tier boxes, no exit ticket.
% The last lesson of Unit 3: it closes the unit's two halves — sampling (3.2–3.3) and
% experiments (3.4) — on one question: what does this study let us conclude, and about whom?
\documentclass[10pt]{article}
\usepackage{apstats-article}
\usepackage{apstats-boxes}

\newcommand{\UnitNumberName}{Unit 3: Collecting Data \quad}
\newcommand{\LessonNumberName}{Lesson 3.5: Inference and Experiments}

\begin{document}
\begin{center}
    {\Large\bfseries \CourseName \\
    {\small\bfseries \UnitNumberName \LessonNumberName}}
\end{center}

\begin{tcolorbox}[colback=sky, colframe=navy, boxrule=0.9pt, arc=2mm,
    left=2mm, right=2mm, top=1.5mm, bottom=1.5mm]
    \textbf{Primary Objective:} Students will read a given simulation of random reshuffles to
    decide whether a difference between treatment groups is statistically significant, use the
    scope of inference to say whether a study supports a cause-and-effect conclusion and to
    whom its result generalizes, and reject a causal claim drawn from a random sample with no
    random assignment.
    \\[2pt]
    \textbf{CED alignment:} Topic 3.7 \quad$\cdot$\quad Big Idea: VAR (Variation and
    Distribution) \quad$\cdot$\quad Skill 4.B \quad$\cdot$\quad LO VAR-3.E \quad$\cdot$\quad
    EK VAR-3.E.1 (inference attributes conclusions to the distribution the data came from),
    VAR-3.E.2 (random assignment lets us call a change too large for chance --- statistically
    significant), VAR-3.E.3 (significant differences between treatment groups are evidence
    the treatments caused the effect), VAR-3.E.4 (random selection of units supports
    generalizing to a larger group)
    \\[2pt]
    \textbf{Lesson model:} \emph{Gradual release --- I do, we do, you do --- all of it inside
    the guided notes.} There is no group activity. The notes name the vocabulary as they teach
    it and deliver the instruction row by row in one \emph{Main Ideas / Notes} table --- each
    idea a definition to complete and a grid of short problems, the first worked by you and the
    rest by the class; the last row, Guided Practice, is worked together with students holding
    the pen; \textbf{the homework is the individual practice}, and students start it in class
    while you circulate. The whole-class debrief is \emph{spoken}.
    \\[2pt]
    \textbf{Compact budget:} the notes are three pages, 14 problems. The simulation is
    \emph{given} as a pre-drawn dotplot: students read it, they never set it up or run it
    (Unit 4 builds simulations; Units 6--9 make significance formal with P-values).
\end{tcolorbox}

\renewcommand{\arraystretch}{1.4}
\boxguard
\begin{skillbox}[Learning Targets \& Key Understandings]{goldbox}
\begin{tabularx}{\linewidth}{@{}X|X@{}}
\textbf{Learning targets (student language)} & \textbf{Key understandings (the ``why'')} \\[2pt]
\hline
\vspace{-6pt}
\begin{itemize}[leftmargin=1.1em, nosep, after=\vspace{2pt}]
  \item I can read a simulation of random reshuffles and decide whether a difference between
        treatment groups is statistically significant.
  \item I can say whether a study supports a cause-and-effect conclusion, and to whom its
        result can be generalized.
  \item I can explain why a random sample, without random assignment, cannot show cause and
        effect.
\end{itemize}
&
\vspace{-6pt}
\begin{itemize}[leftmargin=1.1em, nosep, after=\vspace{2pt}]
  \item Even a useless treatment leaves the groups a little different, from the luck of the
        random assignment. A difference too large to be explained by that luck is
        \emph{statistically significant} (EK VAR-3.E.2), and in a randomized experiment it is
        evidence the treatment caused it (VAR-3.E.3).
  \item Random \emph{assignment} earns the word ``caused''; random \emph{selection} earns
        the population (VAR-3.E.4). They answer different questions.
  \item \textbf{Target misconception:} ``the sample was random, so the study proves cause.''
        A random sample of 500 teens says music-studiers have lower GPAs \emph{in the
        state}; it says nothing about \emph{why}. Only random assignment of the treatment
        rules out confounding.
  \item \textbf{Second misconception:} ``any difference counts.'' A 1.0-word gap is reached
        by 23 of 100 reshuffles --- chance does it all the time.
  \item This lesson closes Unit 3. Unit 4 opens on the same question --- real pattern or
        chance? --- and builds the simulations we only read today.
\end{itemize}
\end{tabularx}
\end{skillbox}

\boxguard
\begin{skillbox}[{Vocabulary, Concepts, \& Theorems --- taught directly in the guided notes}]{greenbox}
\begin{minipage}[t]{0.48\linewidth}
    \begin{itemize}
        \item \textbf{Statistical inference} --- using data to draw conclusions beyond the
              individuals measured (EK VAR-3.E.1)
        \item \textbf{Statistically significant} --- a result too large to be explained by
              chance alone (EK VAR-3.E.2)
        \item \textbf{Simulation of random reshuffles} --- re-randomize the observed results
              many times as if the treatment did nothing; the dotplot shows what chance alone
              produces (read, never run, today)
    \end{itemize}
\end{minipage}
\hfill
\begin{minipage}[t]{0.48\linewidth}
    \begin{itemize}
        \item \textbf{Scope of inference} --- what a study lets us conclude: cause and
              effect, and about whom (the $2\times2$ of random selection $\times$ random
              assignment)
        \item \textbf{Generalize} --- apply the result to the population the individuals
              were randomly selected from (EK VAR-3.E.4)
        \item \textbf{Carried forward} --- random sampling (3.2), confounding and
              observational studies (3.1), random assignment and experimental designs (3.4)
    \end{itemize}
\end{minipage}
\end{skillbox}

% A tabularx never splits, so guard generously ahead of it.
\boxguard[30]
\begin{skillbox}[Lesson at a Glance --- \MeetingLength]{sky}
\renewcommand{\arraystretch}{1.25}
\begin{tabularx}{\linewidth}{@{}l c X X@{}}
\textbf{Phase} & \textbf{Min} & \textbf{Students} & \textbf{Teacher} \\
\hline
Warm-Up & 5 & Selection or assignment, breakfast and grades, and ``could chance alone give
2.5 words?'' & Hand count on item~3; leave ``2.5 words --- chance?'' and the tally on the
board \\
Guided Notes \& Practice & 34 & Fill the vocabulary box and the notes table: problems 1--10 row
by row, then \emph{The Reading App} (11--14) together, pen in their hand & \textbf{I do}
each row's definition lines and first problem (1, 5, 8), \textbf{we do} the rest (24)
$\to$ \textbf{we do} Guided Practice (10) \\
Debrief & 8 & Pens down, papers up --- answer aloud, nothing to fill in & Open on the
almost-right problem~10; put \emph{assigned $\to$ cause, selected $\to$ population} on the
board \\
Close \& Assign & 8 & \textbf{Start the homework in class}, alone and silent & Circulate on
Part~B item~2 for your formative read; preview Lesson 4.1 \\
\end{tabularx}
\end{skillbox}

\boxguard[20]
\begin{skillbox}[Warm-Up --- Activate Prior Knowledge (5 min)]{sky}
\begin{minipage}[t]{0.48\linewidth}
    \textbf{The items and what each seeds}
    \begin{itemize}
        \item \textbf{1.} Random selection or random assignment, and what each is
              \emph{for}. \textbf{Spirals} to 3.2 and 3.4; \textbf{seeds} the Scope of
              Inference row --- selection $\to$ population, assignment $\to$ cause.
        \item \textbf{2.} Breakfast and grades: why an observational study cannot show
              cause. \textbf{Spirals} to 3.1; \textbf{seeds} the crux row.
        \item \textbf{3.} 24 volunteers, silence beats music by 2.5 words: could the luck of
              the assignment do that? \textbf{Seeds} row~1 --- the given simulation answers it
              in problems 2--3.
    \end{itemize}
\end{minipage}
\hfill
\begin{minipage}[t]{0.48\linewidth}
    \textbf{Running it}
    \begin{itemize}
        \item \textbf{Debrief item~1 aloud} and write the two purposes on the board in the
              room's words (``so it's like everybody,'' ``so the groups are fair''). Row~2
              points back at them.
        \item \textbf{Take a hand count on item~3 and do not settle it.} Expect a split:
              ``2.5 is a lot'' against ``anything can happen.'' Both are guesses; the dotplot
              replaces the guess with a count. Write \emph{2.5 words --- chance?} and the tally
              on the board.
        \item Item~2 is a spiral; take 60 seconds and name the confounder.
    \end{itemize}
\end{minipage}
\end{skillbox}

\boxguard
\begin{skillbox}[Guided Notes \& Practice --- I do, then we do (34 min)]{sky}
\textbf{This is the whole lesson.} There is no group activity, and there is no independent
practice set on this page: \textbf{the homework is the individual practice}, started in class.
The notes are one \emph{Main Ideas / Notes} table --- four rows, 14 short problems.
\textbf{In every row you fill the definition lines and work the first problem of each grid; the
class works the rest, pen in their hand.} The vocabulary box is filled \emph{as you go}.
\begin{multicols}{2}
\textbf{Significant (9 min).} Fill the three blanks (individuals / assignment / rarely) and
name \textbf{statistical inference} and \textbf{statistically significant}. Then read the
dotplot \emph{before} anyone writes: one dot is one reshuffle of the 24 real scores into two
new groups, as if music did nothing. Work \textbf{problem~1} yourself. Problems 2 and 3 are
theirs and settle the warm-up vote: 2 of 100. \textbf{Problem~4 is the trap}: ``1.0 is still
more words'' --- make them count the 23 dots at 1.0 or beyond.

\textbf{Scope of Inference (7 min).} Fill the sentence and the $2\times2$ \emph{with} the
class, cell by cell: \emph{was it assigned? was it selected?} Work \textbf{problem~5} (both).
Problems 6 and 7 are theirs; 7 is the volunteers from row~1 --- cause, but only for volunteers
like these.

\textbf{Association or Cause? (8 min) --- the crux.} Fill the three blanks (confounding /
associated / caused). Work \textbf{problem~8}: yes, the association generalizes to all teens in
the state --- the random sample bought that much. Problem~9 is theirs (phone use, study time).
\textbf{Problem~10 is the crux}: take a hand count on ``agree'' \emph{before} anyone explains.
Then fill the callout --- \emph{only random assignment $\to$ cause} --- and make the room say
it back.

\columnbreak
\textbf{We do (about 10 min): The Reading App.} A fresh context, problems 11--14, worked
entirely together. Students hold the pen; you hold the questions: \emph{``How many reshuffles
reached 6?''} \quad \emph{``Who decided who got the app?''} \quad \emph{``Picked from whom?''}
\textbf{Problem~14 is the crux transferred}: a second district's random sample of app users.
It is the one to protect when time is short --- a student who answers both 10 and 14 has
separated \emph{selected} from \emph{assigned}.

\textbf{The pen must actually change hands.} Put problem~11 on them cold, take a hand count on
problem~14 before anyone speaks, and make the room defend the answer. Circulate and demand the
reason:
\begin{itemize}[leftmargin=1.1em, nosep]
  \item \textbf{Problem~11:} ``How often does chance alone get there?''
  \item \textbf{Problem~13:} ``Picked at random \emph{from whom}?''
  \item \textbf{Problem~14:} ``Who decided which students use the app?''
\end{itemize}
While circulating, pick the papers to put up in the debrief --- one \emph{almost}-right
problem~10 (``no, because 500 isn't enough'' or ``no, because it isn't random'') and one that
names the missing random assignment. \textbf{Your formative read comes twice}: here, and again
in the last eight minutes as students start the homework.
\end{multicols}
\end{skillbox}

\boxguard
\begin{skillbox}[Debrief --- whole class (8 min)]{sky}
Pens down, papers up. \textbf{This debrief is spoken} --- nothing on the page to fill in.
\begin{multicols}{2}
\begin{enumerate}[leftmargin=1.3em, nosep, itemsep=3pt]
  \item \textbf{Open on the almost-right problem~10}: ``disagree, because 500 isn't enough''
        or ``because it wasn't random.'' Right verdict, wrong reason --- the sample \emph{was}
        random, and its size is fine. What is missing is random \emph{assignment}. Correct it
        gently and in public.
  \item Write two columns on the board: \emph{Assigned at random? $\to$ caused} and
        \emph{Selected at random? $\to$ the population}. Place the music volunteers, the 500
        teens, and the reading app in them. That board is the lesson.
  \item Take problem~4 back: why is 1.0 word ``not significant'' when 1.0 is bigger than 0?
\end{enumerate}
\columnbreak
\textbf{The sentence to leave on every desk:} \emph{Random selection lets you speak for the
population. Only random assignment lets you say ``caused.''} Make a student say it, not you.

\textbf{If the debrief runs short}, cut item~3 --- item~1 is the only one that cannot be cut.
\textbf{Do not let the debrief eat the last eight minutes} --- the homework start is your
second formative read.
\end{multicols}
\end{skillbox}

\boxguard
\begin{skillbox}[AP Practice --- extra credit, scored 1 to 5]{goldbox}
Two multiple-choice items --- \textbf{MC~1} is the scope of inference for a randomized
experiment on 80 volunteer runners (B: cause, but only for runners like these; (A)
over-generalizes and (C) swaps the two kinds of randomness), and \textbf{MC~2} reads a
simulation in words (46 of 200 reshuffles reach 0.6~kg, so \emph{not} significant; (C) ``proves
no effect'' is the overreach to name) --- then one free-response set on green tea and blood
pressure from a random sample of 1{,}000 adults: \textbf{(a)} the crux in AP clothing (no
cause; name a confounder), \textbf{(b)} generalize the \emph{association} to the state's adults,
and \textbf{(c)}, \textbf{the spiral}, design an experiment (Lesson 3.4). 8 points: 1 per MC,
2 per part.

\textbf{Extra credit, scored 1 to 5, turned in with the homework.} Go over it at the start of
Lesson 4.1 or in the unit review. Hold the AP standard out loud: \emph{in context or it does
not count} --- a part (a) that says ``no, correlation isn't causation'' without saying no one
was assigned to drink tea, and without a confounder tied to both variables, has not earned
both points.
\end{skillbox}

\boxguard
\begin{skillbox}[Homework --- scored 1 to 5, due the first class after two study halls]{goldbox}
Five items in a third context. \textbf{Part~A} (a fruit-first lunch line: 60 students picked at
random, randomly assigned; a given dotplot of 100 reshuffles): \textbf{A1} count the reshuffles
at 0.4 or more (3); \textbf{A2} significant, in context; \textbf{A3} cause \emph{and} to whom
(all 1{,}200 students at the school). \textbf{Part~B} (a city survey of 400 residents picked at
random: walkers report lower stress): \textbf{B1} an AP-style multiple choice (C) with a
required justification; \textbf{B2} the crux restated --- ``chosen at random, so this proves
walking lowers stress.''

\textbf{Start it in the last eight minutes of the period, in class and alone.}
\textbf{Scoring:} 9 points --- A1 1, A2 2, A3 2, B1 1 for the letter and 1 for the
justification, B2 2. \textbf{B2 is the one to read closely}: a student who agrees with the
council member has kept the target misconception into Unit 4. The spiralbox previews Lesson 4.1.

\textbf{Paper homework is the default for this course} --- assign this page unless you have
decided otherwise. \textbf{DeltaMath override (occasional):} on the days you assign DeltaMath
instead, it replaces this page, and this page stays in the packet as extra practice. Never
assign both.
\end{skillbox}

\boxguard
\begin{skillbox}[Watch For (while circulating)]{redbox}
\begin{itemize}[leftmargin=1.2em, nosep]
  \item \textbf{``It was a random sample, so it shows cause''} (notes~10 and 14, HW~B2,
        AP~(a) --- the target misconception). Probe: \emph{``Who decided which teens study with
        music? Then what else could be different about them?''}
  \item \textbf{``Any difference counts''} (notes~4, AP~MC2). Probe: \emph{``Count the dots
        at least that big. Does chance get there often?''}
  \item \textbf{Generalizing from volunteers} (notes~7, AP~MC1). Probe: \emph{``Picked at
        random from whom? Then who can you speak for?''}
  \item \textbf{``Not significant means no effect''} (AP~MC2 (C)). Probe: \emph{``Did the
        simulation show the fertilizer does nothing --- or only that chance could explain
        this result?''}
\end{itemize}
Cold-call prompts: \emph{``Which randomness buys the word `caused'?''} \quad
\emph{``Give me a study that can generalize but cannot conclude cause.''}
\end{skillbox}

\boxguard
\begin{skillbox}[Close \& Assign (8 min)]{goldbox}
\textbf{Students start the homework here, in class, alone and silent --- this is the individual
practice, and the first minutes of it are your best formative read.} Hand the launch first ---
five items on paper, scored, due the first class after two study halls --- and point out that
the AP Practice page is extra credit, turned in with the homework. Circulate on \textbf{Part~B
item~2}. Sort what you see into three piles as you walk --- disagreed and said no one was
assigned to walk, with a confounder (ready); disagreed with ``correlation isn't causation'' and
nothing more (ready, needs the language); agreed with the council member (the misconception
survived) --- and open the Unit 3 review by returning to the two-column board if that third
pile ran past a third of the room. Then close on what changed: \emph{``You walked in asking
whether 2.5 words could be luck. You walk out knowing which randomness lets you say
`caused.'\,''} \textbf{Preview:} Lesson 4.1, \emph{Random and Non-Random Patterns?} --- Unit 4
starts with the question the dotplot answered today: is a pattern real, or just chance?
\end{skillbox}

% ── Teacher notes — one per component, in packet order (four: there is no activity) ──
\begin{teachernote}[Warm-Up]
Item~3 is the seed and item~1 the quiet setup. Give five minutes total and spend at most 60
seconds on item~2. On item~3 do not settle the vote --- write \emph{2.5 words --- chance?} and
the tally on the board, and say the notes count it. On item~1, the purposes in the room's
words are what row~2 points back at; write them down.
\end{teachernote}

\begin{teachernote}[Guided Notes \& Practice]
\textbf{Pace it 24 / 10, then 8 for the debrief, then 8 for the homework start.} Four rows, 14
problems. \textbf{You work 1, 5, and 8}; the class works the rest. The dotplot is
\emph{given}: explain once what one dot is and move on --- do not build or run the simulation
(that is Unit 4's job). Scope of Inference must move, seven minutes, because the time belongs
to Association or Cause?, which carries the crux. \textbf{Problem~4 is the trap} (1.0 word
reads as ``a difference''); \textbf{problem~10 is the crux} (a random sample read as proof of
cause). Take a hand count on 10 before any explanation; the callout sentence is the
resolution. \textbf{The Reading App is where the pen changes hands}; problem~14 is the crux
transferred and the one to protect. If you only get through 11, 12, and 14, that is the right
trade. \textbf{Your second formative read is the homework start}: Part~B item~2, three piles.
\textbf{Debrief (8 min), spoken}, opening on the almost-right problem~10. \textbf{Do not let the
debrief eat the last eight minutes} --- the homework start is not optional padding.
\end{teachernote}

\begin{teachernote}[AP Practice]
\textbf{This page is extra credit, scored 1 to 5} --- the graded practice is the homework
behind it. \textbf{FRQ part (a) is the one worth reading closely}: ``no'' alone, or ``no,
correlation is not causation,'' earns at most half; full credit says no one was assigned to
drink tea and names a confounder linked to both tea and blood pressure. On MC~1, (C) is the
swap of the two randomnesses --- if more than a handful choose it, return to the two-column
board. Part (c) is a Lesson 3.4 spiral: random assignment to tea or a comparison drink, and a
named response.
\end{teachernote}

\begin{teachernote}[Homework]
\textbf{Scored 1 to 5 (9 points), due the first class after two study halls.} The eight
minutes at the end of the period are a \emph{start}: put students on A1 and B2 first while you
can still catch a wrong start. \textbf{B2 is the item to read first when scoring} --- a student
who sides with the council member has kept the target misconception. On A2, ``yes'' with no
reference to the 3 of 100 reshuffles earns 1 of 2. On A3, both halves are needed: cause (random
assignment) and the school's 1{,}200 students (random selection) --- ``all students'' is an
over-generalization. On B1, (C) with no justification earns 1 of 2. \textbf{Paper is the
default.} On the occasions you assign DeltaMath in its place, tell students this page is now
optional practice. Do not assign both.
\end{teachernote}
\end{document}
